\begin{statementbox}
  \textbf{\textcolor{CStatement}{先解决什么问题？}}

  已知 $y=f(x)$，如果把\textbf{整条图像}向左或向右挪几格，怎样不重新画图，直接写出新函数？

  \medskip
  \textbf{\textcolor{CStatement}{核心结论：左加右减}}

  设平移距离 $a>0$。把原函数中的\textbf{每一个 $x$} 换掉，就能得到新解析式：
  \[
    \boxed{\text{向左移 }a\text{ 格：}\;y=f(x+a)}
    \qquad
    \boxed{\text{向右移 }a\text{ 格：}\;y=f(x-a)}
  \]
  注意：\textbf{图像往哪边走，与式子中加减号的方向相反。}
  对应点的纵坐标保持不变：$(u,f(u))$ 左移后为 $(u-a,f(u))$，右移后为 $(u+a,f(u))$。

  \medskip
  \textbf{\textcolor{CStatement}{看图记住：顶点跟着走}}
  \begin{center}
  \begin{tikzpicture}[x=0.94cm,y=0.77cm,>=stealth,font=\small]
    \draw[->,gray!75] (-3.9,0)--(3.95,0) node[right] {$x$};
    \draw[->,gray!75] (0,-0.3)--(0,2.7) node[above left,xshift=-20pt] {$y$};
    \draw[very thick,blue!80!black,domain=-1.49:1.49,samples=55] plot (\x,{\x*\x});
    \draw[very thick,teal!75!black,domain=-3.49:-0.51,samples=55] plot (\x,{(\x+2)*(\x+2)});
    \draw[very thick,orange!85!black,domain=0.51:3.49,samples=55] plot (\x,{(\x-2)*(\x-2)});
    \fill[teal!75!black] (-2,0) circle (1.9pt);
    \fill[blue!80!black] (0,0) circle (1.9pt);
    \fill[orange!85!black] (2,0) circle (1.9pt);
    \node[below,teal!75!black] at (-2,-0.04) {$-2$};
    \node[below,blue!80!black] at (0,-0.04) {$0$};
    \node[below,orange!85!black] at (2,-0.04) {$2$};
    \node[teal!75!black,above] at (-2.55,2.3) {$y=(x+2)^2$};
    \node[blue!80!black,above] at (0,2.36) {$y=x^2$};
    \node[orange!85!black,above] at (2.5,2.3) {$y=(x-2)^2$};
    \draw[->,teal!75!black,thick] (-0.24,-0.72)--(-1.76,-0.72)
      node[midway,below] {左移 $2$};
    \draw[->,orange!85!black,thick] (0.24,-0.72)--(1.76,-0.72)
      node[midway,below] {右移 $2$};
  \end{tikzpicture}

  \end{center}

  \textbf{\textcolor{CStatement}{适用范围与定义域}}

  任何函数的整条图像都可以左右平移，不需要它的定义域事先“容得下”平移。
  如果原来只能取 $x\in D$，左移后要求 $x+a\in D$，右移后要求 $x-a\in D$；\textbf{定义域也随图像一起移动}。
  例如，原定义域为 $[1,4]$，左移 $2$ 后为 $[-1,2]$，右移 $2$ 后为 $[3,6]$。

  \textbf{统一写法：} $y=f(x-h)$。$h>0$ 向右移 $h$ 格，$h<0$ 向左移 $|h|$ 格，$h=0$ 不动。
\end{statementbox}
